\scanpage{441}
\mspage{353}
Par ailleurs $f(x)\in\mathrm{Ob}\,\underline T$ (puisque $x\in\underline{\mathrm{Ob}}\,\underline S$, \unsure{$\therefore$|et} $\underline S\overset{\mathrm{def}}{=}f^{-1}(\underline T)$). Donc on peut appliquer 4.13.22 $a\Longrightarrow c$ à la situation $(N,\pi_0,\pi,\underline S)$ et au système inductif $(f(x_i))_{i\in I}$, pour conclure que la sous-catégorie $I'$ de $I$, formée des $i$ tels que $f(x_i)\in\deleted{\unsure{\mathrm{Ob}}}\underline S\deleted{\unsure{S}}$ i.e \unsure{alors|tels} que $x_i\in\underline S$ donc $x_i\in\underline S(\pi)$, est \uline{cofinale} dans $I$. Donc on a
\[
x\simeq\underset{i\in I'}{\underrightarrow{\lim}}\,x_{i\unsure{\deleted{{}'}|{}'|\ast}},
\]
avec cette fois les $x_{i\deleted{{}'}}$ dans $\underline S(\pi)$, et $I'$ grde devant $\pi$ (car cofinale \inserted{et pleine} dans $I$, qui est grde devant $\pi$), ce qui prouve que $x$ est ds l'image essentielle de $\mathrm{Ind}_\pi(\underline S(\pi))$, qed.

\uline{Questions} 4.13.33. On vient de voir que si $f:M\unsure{\text{---}|\longrightarrow}N$ est un foncteur accessible entre catégories accessibles, alors l'image inverse d'une \deleted{\unsure{sous-catég}} partie accessible de $\mathrm{Ob}\,N$ est une partie accessible de $\mathrm{Ob}\,M$.
% slanted
\margin{Faux, $=$ si on prend la \unsure{$G$-clôture|$\mathcal G$-clôture} de l'image essentielle (``image $G$-essentielle'')}
Je conjecture que l'énoncé \unsure{similaire|symétrique}, \inserted{non pas} pour les images directes, mais pour les \uline{images essentielles}, est
% archive sheet number: 440
